% BOOK-FILE-ID: f-d6b88598ed80
% subsection: 隣接3項間漸化式型 \, $\displaystyle a_{n+2}+pa_{n+1}+qa_n=0$

\subsection{\texorpdfstring{隣接3項型  $a_{n+2}+pa_{n+1}+qa_n=0$}{隣接3項型}}\label{節：隣接3項型}
\subsubsection{\texorpdfstring{付加項がない場合  $a_{n+2}+pa_{n+1}+qa_n=0$}{付加項がない場合}}\label{節：付加項がない場合}

今までは隣接2項間の漸化式について考えてきたが,過去2項の影響を受けるような漸化式を考えてみよう.
この漸化式の有名な例にフィボナッチ数列の漸化式などもある.

一般的な形は
\[
  a_{n+2}+pa_{n+1}+qa_n=0
\]
で表される.
いままで数々の2項間漸化式を扱ってきた.
したがって,今回は3項間をなんとかして2項間にできないか挑戦してみよう.

過去の天才らのアイデアを自力で出すのは難しいので今回も結論から示す.
その手法は
\BookMathLayout{\[
  \begin{aligned}
    a_{n+2}+pa_{n+1}+qa_n&=0\\
    &\iff a_{n+2}-\alpha a_{n+1}\\
    &\qquad=\beta(a_{n+1}-\alpha a_n)
  \end{aligned}
\]}{\[
  \begin{aligned}
    a_{n+2}+pa_{n+1}+qa_n&=0\\
    &\iff a_{n+2}-\alpha a_{n+1}\\
    &\qquad=\beta(a_{n+1}-\alpha a_n)
  \end{aligned}
\]}
となる解$\alpha,\beta$を考えることで,\,$a_n$に関しての3項間の漸化式が$a_n$の階差に関する2項間の漸化式になるのだ.
$a_n$の階差に関する2項間漸化式を式変形して係数比較すると,
\begin{align*}
  a_{n+2}-\alpha a_{n+1}&=\beta(a_{n+1}-\alpha a_{n})\\
  a_{n+2}-(\alpha+\beta)a_{n+1}&=-\alpha\beta a_n
\end{align*}
より,
\[
  \alpha+\beta=-p,\qquad\alpha\beta=q
\]
解と係数の関係より,\,$\alpha,\beta$は$x^2+px+q=0$の二根である.
この方程式もまた\bookterm{characteristic}{特性方程式}という.
すなわち,形式的には$a_{n+2}$を$x^2$,\,$a_{n+1}$を$x$,\,$a_n$を1に置き換えた二次方程式を解くことで等比型に帰着させられるのである.

なお$q=0$なら先に元の式を使う.\,$a_{n+2}=-pa_{n+1}$なので,\,$p\ne0$では
\[
  a_n=a_2(-p)^{n-2}\quad(n\ge2)
\]
となり,\,$a_1$は与えられた値を保つ.\,$p=0$でもあるなら$n\ge3$で$a_n=0$である.
以下では$q\ne0$,すなわち二根とも0でないとする.

$\alpha\ne\beta$なら,
\[
  \begin{aligned}
    b_n&=a_{n+1}-\alpha a_n,\\
    c_n&=a_{n+1}-\beta a_n
  \end{aligned}
\]
はそれぞれ公比$\beta,\alpha$の\bookterm{geometric}{等比数列}になる.
\begin{align*}
  b_{n+1}&=\beta b_n,& b_1&=a_2-\alpha a_1,\\
  c_{n+1}&=\alpha c_n,& c_1&=a_2-\beta a_1.
\end{align*}
求まったのは元の数列ではなく,2項をまとめた量である.図の上段は元の2項の組,下段は$b_n$だけの漸化式を表す.
\BookFigThreeTerm

$b_n-c_n=(\beta-\alpha)a_n$だから,
\BookMathLayout{\[
    a_n=\frac{(a_2-\alpha a_1)\beta^{n-1}-(a_2-\beta a_1)\alpha^{n-1}}{\beta-\alpha}.
  \]}{\[
    \begin{aligned}
      a_n&=\frac{a_2-\alpha a_1}{\beta-\alpha}\beta^{n-1}\\
      &\quad-\frac{a_2-\beta a_1}{\beta-\alpha}\alpha^{n-1}.
    \end{aligned}
  \]}
係数は$n$によらないので,\,$a_n=A\alpha^{n-1}+B\beta^{n-1}$と書ける.
重解では作れた一つの等比数列から,\,$a_{n+1}-\alpha a_n$の値は分かる.
そこで,今度はこの関係を$a_n$について考えればよい.
$b_n=a_{n+1}-\alpha a_n$を等比型として解き,残る式から復元する.例題(2)で確かめよう.

\begin{bookpoint}
  \textbf{POINT}\par
  $a_{n+2}$を$x^2$,\,$a_{n+1}$を$x$,\,$a_n$を1に置き換えた2次方程式を解き,得られた2解$\alpha,\beta$をもって立てられる2漸化式
  \begin{align*}
    \begin{cases}
      a_{n+2}-\alpha a_{n+1}=\beta(a_{n+1}-\alpha a_n) &\\
      a_{n+2}-\beta a_{n+1}=\alpha(a_{n+1}-\beta a_n) &\\
    \end{cases}
  \end{align*}
  を解けば良い.
\end{bookpoint}
\noindent \bookblocklink{example-three-term-1}{problem}{answer}{【例題】}\par\nobreak
\begin{enumerate}[label=(\arabic*),parsep=3pt,format=\bookitemlink{example-three-term-1}{problem}{answer}]
  \item $a_1=1,\,a_2=1,\, a_{n+2}=5a_{n+1}-6a_n$
  \item $a_1=1,\,a_2=4,\, a_{n+2}=4a_{n+1}-4a_n$
  \item $a_1=2,\,a_2=1,\, a_{n+2}=a_{n+1}+6a_n$
\end{enumerate}

\noindent\bookblocklink{example-three-term-1}{hint}{problem}{【方針】}\par\nobreak
\begin{enumerate}[label=(\arabic*),parsep=3pt,format=\bookitemlink{example-three-term-1}{hint}{problem}]
  \item \bookterm{characteristic}{特性方程式}$x^2-5x+6=0$の解は異なる2つの実数解です.それぞれの根で二項をまとめ,\,2本の\bookterm{geometric}{等比数列}を作りましょう.
  \item \bookterm{characteristic}{特性方程式}$x^2-4x+4=0$を解くと重解になります.\bookterm{geometric}{等比数列}は1本しか作れないことに注意しましょう.
  \item \bookterm{characteristic}{特性方程式}の解の1つが負になるパターンです.符号に注意して計算しましょう.
\end{enumerate}

\noindent\bookblocklink{example-three-term-1}{answer}{problem}{【解答】}\par\nobreak
\begin{enumerate}[label=(\arabic*),parsep=3pt,format=\bookitemlink{example-three-term-1}{answer}{problem}]
  \item $a_n=2^n-3^{n-1}$
  \item $a_n=n\cdot2^{n-1}$
  \item $a_n=(-2)^{n-1}+3^{n-1}$
\end{enumerate}

\noindent\bookblocklink{example-three-term-1}{solution}{problem}{【解法】}\par\nobreak
\begin{enumerate}[label=(\arabic*),parsep=3pt,format=\bookitemlink{example-three-term-1}{solution}{problem}]
  \item \bookterm{characteristic}{特性方程式}$x^2-5x+6=(x-2)(x-3)=0$の根は$2,3$である.
    \begin{align*}
      a_{n+2}-2a_{n+1}&=3(a_{n+1}-2a_n),\\
      a_{n+2}-3a_{n+1}&=2(a_{n+1}-3a_n).
    \end{align*}
    $b_n=a_{n+1}-2a_n,\ c_n=a_{n+1}-3a_n$と置くと,
    \[
      \begin{aligned}
        b_1&=-1,& b_n&=-3^{n-1},\\
        c_1&=-2,& c_n&=-2^n.
      \end{aligned}
    \]
    $b_n-c_n=a_n$より,\,$a_n=2^n-3^{n-1}$を得る.
  \item \bookterm{characteristic}{特性方程式}$x^2-4x+4=(x-2)^2=0$は重解$2$を持つ.
    $b_n=a_{n+1}-2a_n$と置くと,
    \[
      \begin{aligned}
        b_{n+1}&=2(a_{n+1}-2a_n)=2b_n,\\
        b_1&=4-2=2.
      \end{aligned}
    \]
    従って$b_n=2^n$であり,\,$a_{n+1}-2a_n=2^n$となる.
    $2^{n+1}$で割り,\,$c_n=\frac{a_n}{2^n}$と置けば,
    \[
      c_{n+1}-c_n=\frac12,\qquad c_1=\frac12.
    \]
    差を足すと$c_n=\frac{n}{2}$なので,\,$a_n=2^nc_n=n2^{n-1}$となる.
  \item \bookterm{characteristic}{特性方程式}$x^2-x-6=(x-3)(x+2)=0$の根は$3,-2$である.
    \begin{align*}
      a_{n+2}-3a_{n+1}&=-2(a_{n+1}-3a_n),\\
      a_{n+2}+2a_{n+1}&=3(a_{n+1}+2a_n).
    \end{align*}
    $b_n=a_{n+1}-3a_n,\ c_n=a_{n+1}+2a_n$と置くと,
    \[
      \begin{aligned}
        b_1&=-5,& b_n&=-5(-2)^{n-1},\\
        c_1&=5,& c_n&=5\cdot3^{n-1}.
      \end{aligned}
    \]
    $c_n-b_n=5a_n$より,\,$a_n=(-2)^{n-1}+3^{n-1}$を得る.
\end{enumerate}

\paragraphbf{2つの解の和で表せる理由}\label{解釈：三項間と解全体}
$a_n=\alpha^{n-1}$と$a_n=\beta^{n-1}$は,どちらも$a_{n+2}+pa_{n+1}+qa_n=0$を満たす.
この漸化式は,解を定数倍して足しても右辺が0のまま変わらないので,\,$a_n=A\alpha^{n-1}+B\beta^{n-1}$も漸化式を満たす.
さらに$\alpha\ne\beta$なら,初めの2項についての連立方程式$A+B=a_1,\ A\alpha+B\beta=a_2$から$A,B$がただ1組に定まる.
したがって,どのような\bookterm{initial}{初期条件}を与えても,解はこの形で表せる.この背景は付録の\chapterref{章：解法の背景と広がり}で扱う.

\paragraphbf{重解で一次式が残る理由}
例題(2)で$n$が掛かった理由を,一般の0でない重解$\alpha$でも確かめよう.
$b_n=a_{n+1}-\alpha a_n$は公比$\alpha$の等比型なので,
\[
  a_{n+1}-\alpha a_n=(a_2-\alpha a_1)\alpha^{n-1}
\]
である.ここで$c_n=\frac{a_n}{\alpha^{n-1}}$と置く.
両辺を$\alpha^n$で割ると,
\[
  c_{n+1}-c_n=\frac{a_2-\alpha a_1}{\alpha}
\]
となる.つまり,\,$c_n$は公差$\frac{a_2-\alpha a_1}{\alpha}$の等差数列である.
$c_1=a_1$から$c_n$を求め,\,$a_n=\alpha^{n-1}c_n$に戻すと,
\BookMathLayout{\[
    a_n=\left\{a_1+(n-1)\left(\frac{a_2}{\alpha}-a_1\right)\right\}\alpha^{n-1}
  \]}{\[
    \begin{aligned}
      a_n&=\left\{a_1+(n-1)\left(\frac{a_2}{\alpha}-a_1\right)\right\}\\
      &\qquad\cdot\alpha^{n-1}
    \end{aligned}
  \]}
を得る.この一定の差を$n-1$回足したことが,一次式を作ったのである.\footnote{$\alpha^{n-1}$で割ってから階差を取るこの方法は,隣接4項型以上で特性方程式の解が3重,4重に重なる場合にも使える.付録の\sectionref*{節：差分で線形漸化式を解く}では,これを別解として一般の場合に広げている.}
$\alpha^{n-1}$と$(n-1)\alpha^{n-1}$の初めの2項は$(1,\alpha)$と$(0,\alpha)$である.
$\alpha\ne0$なので,どのような$a_1,a_2$に対しても係数が1組に定まり,重解の場合もすべての解をこの形で表せる.

\paragraphbf{別解：隣り合う項の比を使う}

ここで扱う別解は,特に計算において有利になるようなものでも,汎用性があるようなものでもない.
しかしながら他の解法とのつながりという観点で大きな意味を持つため,解説する.

隣接3項間漸化式
\[
  a_{n+2}+pa_{n+1}+qa_n=0
\]
について,すべての項が0でない場合を考える.
両辺を$a_{n+1}$で割って,
\[
  \frac{a_{n+2}}{a_{n+1}}+p+q\frac{a_n}{a_{n+1}}=0
\]
ここで$b_n=\frac{a_{n+1}}{a_n}$と置いて式を整理すると,
\[
  b_{n+1}=-p-\frac{q}{b_n}
\]
となり,一次分数型になる.初項は$b_1=\frac{a_2}{a_1}$であり,この式を解いた後は
\[
  a_n=a_1\prod_{k=1}^{n-1}b_k
\]
で元の項に戻せる.\,$n=1$の空積は$1$とする.
行列を用いた詳しい説明は\sectionref*{節：1つの行列で見渡す}で扱う.

\subsubsection{\texorpdfstring{定数の付加項がある場合  $a_{n+2}+pa_{n+1}+qa_n=r$}{定数の付加項がある場合}}\label{節：定数の付加項がある場合}
定数の付加項$r$も,\bookterm{fixed}{不動点}となる定数解$\alpha$を引いて消す.
$b_n=a_n-\alpha$が$b_{n+2}+pb_{n+1}+qb_n=0$を満たす条件は,
\[
  r=(1+p+q)\alpha.
\]
$1+p+q\ne0$なら$\alpha=\frac{r}{1+p+q}$である.
$b_1=a_1-\alpha,\ b_2=a_2-\alpha$から前節の隣接3項型を解き,\,$a_n=b_n+\alpha$へ戻す.
$1+p+q=0$の場合は,例題の後で階差を使って考える.

\noindent \bookblocklink{example-three-term-2}{problem}{answer}{【例題】}\par\nobreak
\begin{enumerate}[label=(\arabic*),parsep=3pt,format=\bookitemlink{example-three-term-2}{problem}{answer}]
  \item $a_1=2,\,a_2=2,\, a_{n+2}=5a_{n+1}-6a_n+2$
  \item $a_1=1,\,a_2=2,\, a_{n+2}=4a_{n+1}-4a_n+3$
\end{enumerate}

\noindent\bookblocklink{example-three-term-2}{hint}{problem}{【方針】}\par\nobreak
\begin{enumerate}[label=(\arabic*),parsep=3pt,format=\bookitemlink{example-three-term-2}{hint}{problem}]
  \item まず\bookterm{fixed}{不動点}$\alpha$を求め,\,$b_n=a_n-\alpha$と置いて前節の二次\bookterm{characteristic}{特性方程式}型に帰着させましょう.
  \item こちらは重解のパターンです.同じ手順で$\alpha$を求めてから前節の重解の場合の処理を行いましょう.
\end{enumerate}

\noindent\bookblocklink{example-three-term-2}{answer}{problem}{【解答】}\par\nobreak
\begin{enumerate}[label=(\arabic*),parsep=3pt,format=\bookitemlink{example-three-term-2}{answer}{problem}]
  \item $a_n=1+2^n-3^{n-1}$
  \item $a_n=3+2^{n-2}(3n-7)$
\end{enumerate}

\noindent\bookblocklink{example-three-term-2}{solution}{problem}{【解法】}\par\nobreak
\begin{enumerate}[label=(\arabic*),parsep=3pt,format=\bookitemlink{example-three-term-2}{solution}{problem}]
  \item $r=2,\ p=-5,\ q=6$より\BookMathLayout{}{\\}%
    $\alpha=\frac{2}{1-5+6}=1$である.
    $b_n=a_n-1$と置くと,
    \[
      b_{n+2}=5b_{n+1}-6b_n,\qquad b_1=b_2=1.
    \]
    \bookproblemref{example-three-term-1}{1}{前の例題(1)}と同じ漸化式・初期条件だから,\,$b_n=2^n-3^{n-1}$となる.
    従って$a_n=1+2^n-3^{n-1}$である.
  \item $r=3,\ p=-4,\ q=4$より$\alpha=3$である.\,$b_n=a_n-3$と置くと,
    \[
      b_{n+2}=4b_{n+1}-4b_n,\qquad b_1=-2,\quad b_2=-1.
    \]
    特性方程式は重解$2$を持つ.\,$c_n=b_{n+1}-2b_n$なら$c_{n+1}=2c_n,\ c_1=3$なので,\,$c_n=3\cdot2^{n-1}$である.
    $b_{n+1}-2b_n=c_n$を$2^{n+1}$で割り,\,$f_n=\frac{b_n}{2^n}$と置くと,
    \[
      f_{n+1}-f_n=\frac34,\qquad f_1=-1.
    \]
    従って$f_n=-1+\frac34(n-1)=\frac{3n-7}{4}$であり,
    \[
      a_n=3+b_n=3+2^{n-2}(3n-7).
    \]
\end{enumerate}


\paragraphbf{定数でずれを戻せないとき}

$p+q+1=0$かつ$r\ne0$のときは,\,$r=(p+q+1)\alpha$を満たす定数$\alpha$が存在しない.
一方,\,$r=0$なら任意の定数が条件を満たし,元の式はすでに付加項のない形である.

定数解が存在しない場合は,隣り合う項の差を$d_n=a_{n+1}-a_n$と置く.
$-(p+1)=q$を使うと,
\[
  \begin{aligned}
    d_{n+1}&=a_{n+2}-a_{n+1}\\
    &=-(p+1)a_{n+1}-qa_n+r\\
    &=q(a_{n+1}-a_n)+r=qd_n+r
  \end{aligned}
\]
となり,特性方程式型または\bookterm{difference}{階差}型に帰着できる.
$d_1=a_2-a_1$から解き,\,$a_n=a_1+\sum_{k=1}^{n-1}d_k$として求める.
